
\documentclass[12pt]{article}
\usepackage[margin=0.75in]{geometry}
\usepackage{fontspec}
\setmainfont{Latin Modern Roman}
\usepackage{microtype}
\usepackage{multicol}
\usepackage{fancyhdr}
\usepackage{titlesec}
\usepackage{enumitem}
\usepackage{xcolor}
\usepackage[hidelinks]{hyperref}
\setlength{\columnsep}{0.28in}
\setlength{\parindent}{1.1em}
\setlength{\parskip}{0.35em}
\setlength{\headheight}{15pt}
\titleformat{\section}{\large\bfseries\scshape}{}{0pt}{}
\titleformat{\subsection}{\normalsize\bfseries}{}{0pt}{}
\pagestyle{fancy}
\fancyhf{}
\fancyfoot[C]{\thepage}
\renewcommand{\headrulewidth}{0.4pt}
\newcommand{\focusbox}[1]{\par\vspace{0.4em}\noindent\begingroup\setlength{\fboxsep}{6pt}\colorbox{gray!12}{\begin{minipage}{\dimexpr\linewidth-2\fboxsep\relax}#1\end{minipage}}\endgroup\par\vspace{0.5em}}
\newcommand{\studentline}{\vspace{0.35em}\hrule\vspace{0.65em}}
\newcommand{\shortline}{\vspace{0.25em}\hrule\vspace{0.45em}}

\fancyhead[L]{Whately: Argument Lab}
\fancyhead[R]{Student Handout}
\begin{document}
\begin{center}
{\LARGE\bfseries Week 1 Argument Lab}\par\vspace{0.25em}
{\large\scshape What Follows From What?}\par\vspace{0.6em}
\rule{0.82\textwidth}{0.4pt}
\end{center}

\section*{Name: \rule{2.4in}{0.4pt} \hfill Date: \rule{1.3in}{0.4pt}}

\section*{Purpose}
Whately says logic studies reasoning. Today you will practice finding the structure of an argument: the conclusion, the stated reason, the hidden assumption, and the question of whether the conclusion actually follows. The first item repeats a pattern from the reading. The rest are new arguments.

\section*{Part I: One Model From the Reading}
Use this example to remember the four steps.

\subsection*{Example}
\textbf{Argument:} "This rule is popular with most students, so it must be fair."

\begin{itemize}[leftmargin=*]
\item \textbf{Conclusion:} The rule must be fair.
\item \textbf{Stated reason:} Most students like it.
\item \textbf{Hidden assumption:} Whatever most students like is fair.
\item \textbf{Test:} The conclusion does not necessarily follow. Popularity and fairness are not the same thing.
\end{itemize}

\section*{Part II: Fresh Arguments}
For each argument, name the parts and test the reasoning. Some arguments are weak. At least one is stronger than it first appears. Mark one response with a star as a possible portfolio artifact.

\subsection*{1. Argument}
"The city placed a bronze statue of this man in the town square, so he must have lived nobly."

\textbf{Conclusion:}\studentline
\textbf{Stated reason:}\studentline
\textbf{Hidden assumption:}\studentline
\textbf{Does the conclusion follow? Why or why not?}\studentline\studentline

\subsection*{2. Argument}
"The first essay uses more difficult words than the second essay, so the first essay must have better reasoning."

\textbf{Conclusion:}\studentline
\textbf{Stated reason:}\studentline
\textbf{Hidden assumption:}\studentline
\textbf{Does the conclusion follow? Why or why not?}\studentline\studentline

\subsection*{3. Argument}
"Our class has always done the assignment this way, so this way must be the best way to do it."

\textbf{Conclusion:}\studentline
\textbf{Stated reason:}\studentline
\textbf{Hidden assumption:}\studentline
\textbf{Does the conclusion follow? Why or why not?}\studentline\studentline

\subsection*{4. Argument}
"Every just rule should treat the same action the same way. This rule punishes one student and excuses another student for the same action. Therefore, this rule is not just."

\textbf{Conclusion:}\studentline
\textbf{Stated reason:}\studentline
\textbf{Hidden assumption:}\studentline
\textbf{Does the conclusion follow? Why or why not?}\studentline\studentline

\subsection*{5. Literary Anchor: Macbeth or Caesar}
Choose one claim below and analyze it with the four-step test.

\textbf{Macbeth claim:} "The witches predicted Macbeth's future, so Macbeth is not responsible for what he later chooses."

\textbf{Caesar claim:} "Caesar may become dangerous later, so Brutus is justified in killing him now."

\textbf{Conclusion:}\studentline
\textbf{Stated reason:}\studentline
\textbf{Hidden assumption:}\studentline
\textbf{Does the conclusion follow? Why or why not?}\studentline\studentline

\section*{Part III: Repair the Reasoning}
Choose \textbf{one weak argument} from Part II. Rewrite it so that the conclusion is either better supported or more careful. You may change the conclusion, add a missing reason, or narrow the claim.

\focusbox{Weak argument number: \rule{0.8in}{0.4pt}\hfill My repaired version:}
\vspace{1.35in}

\section*{Part IV: Create Your Own Test Case}
Write one original argument about a school rule, a speaker in \textit{Macbeth} or \textit{Julius Caesar}, a public speech, a sports decision, or a family disagreement. Then analyze your own argument.

\textbf{My argument:}\studentline\studentline
\textbf{Conclusion:}\shortline
\textbf{Stated reason:}\shortline
\textbf{Hidden assumption:}\shortline
\textbf{Does the conclusion follow? Why or why not?}\studentline

\section*{Exit Claim}
Complete this claim in one strong paragraph:

\focusbox{\textbf{Portfolio exit claim:} Logic helps a reader most when \rule{2.0in}{0.4pt} because \rule{2.2in}{0.4pt}. Use one example from Whately, \textit{Macbeth}, or \textit{Julius Caesar}.}

\vspace{1.2in}
\end{document}
